\chapter{Complete Examples}
\label{chap:complete-examples}

The examples in this chapter combine the individual keyword pages into complete modelling workflows. Each section describes one physical problem and presents the native FEMaster and supported Abaqus representation side by side. The examples intentionally keep explicit regions, properties, boundary conditions, loads, and analysis controls so that the translation between both dialects remains visible.

The two readers do not expose exactly the same abstractions. Native FEMaster defines reusable support and load collectors at model level and selects them in \cmd{LOADCASE}. Abaqus input places supported \cmd{BOUNDARY} and load cards in one \cmd{STEP}. Native truss properties use \cmd{TRUSS SECTION}; the current Abaqus reader maps a pure \val{T3D2} set supplied through \cmd{SOLID SECTION} to the same internal TrussSection. These examples always use the syntax actually implemented by the readers rather than syntax merely available in the full Abaqus product.

\section{Linear static truss}

The first model is a two-node axial bar in a consistent N--mm unit system. Node 1 is fixed and node 2 is guided so only axial translation remains free. A 1000 N force acts in the global \(x\)-direction. With \(E=210000\) MPa, \(L=1000\) mm, and \(A=100\) mm\(^2\), the analytical displacement is
\[
u_x=\frac{FL}{EA}\approx0.047619\ \mathrm{mm}.
\]
The case is therefore useful both as a syntax template and as a basic verification model.

\begin{dialectexamples}
\begin{femasterdeck}
*HEADING
Linear static axial truss

*NODE
1,    0., 0., 0.
2, 1000., 0., 0.
*ELEMENT, TYPE=T3, ELSET=BAR
1, 1, 2
*NSET, NAME=FIXED
1
*NSET, NAME=TIP
2

*MATERIAL, NAME=STEEL
*ELASTIC
210000., 0.3
*TRUSS SECTION, ELSET=BAR,
 MATERIAL=STEEL
100.

*SUPPORT, SUPPORT_COLLECTOR=BC
FIXED, 0., 0., 0., NAN, NAN, NAN
TIP, NAN, 0., 0., NAN, NAN, NAN
*CLOAD, LOAD_COLLECTOR=AXIAL
TIP, 1000., 0., 0., 0., 0., 0.

*LOADCASE, TYPE=LINEARSTATIC
*SUPPORTS
BC
*LOADS
AXIAL
*SOLVER, DEVICE=CPU, METHOD=DIRECT
*CONSTRAINTMETHOD, TYPE=NULLSPACE
*END
\end{femasterdeck}
\begin{abaqusdeck}
*HEADING
Linear static axial truss

*NODE
1,    0., 0., 0.
2, 1000., 0., 0.
*ELEMENT, TYPE=T3D2, ELSET=BAR
1, 1, 2
*NSET, NSET=FIXED
1
*NSET, NSET=TIP
2

*MATERIAL, NAME=STEEL
*ELASTIC
210000., 0.3
*SOLID SECTION, ELSET=BAR,
 MATERIAL=STEEL
100.

*STEP, NAME=AXIAL_TENSION
*STATIC
*BOUNDARY
FIXED, 1, 3, 0.
TIP, 2, 3, 0.
*CLOAD
TIP, 1, 1000.
*END STEP
\end{abaqusdeck}
\end{dialectexamples}

\section{Reusable Part and translated Instance}

This example demonstrates the optional Part/Assembly hierarchy. The material is a root-level named resource. The reusable truss topology and its Section assignment live inside the Part. One Instance translates the complete Part by 500 mm in \(x\). Assembly-level node sets then resolve the local node labels through \key{INSTANCE}. The source Part remains unchanged; transformed coordinates are created only when the model is compiled.

\begin{dialectexamples}
\begin{femasterdeck}
*MATERIAL, NAME=STEEL
*ELASTIC
210000., 0.3

*PART, NAME=ROD
*NODE
1,   0., 0., 0.
2, 500., 0., 0.
*ELEMENT, TYPE=T3, ELSET=BAR
1, 1, 2
*TRUSS SECTION, ELSET=BAR,
 MATERIAL=STEEL
50.
*END PART

*ASSEMBLY
*INSTANCE, NAME=ROD_A, PART=ROD
500., 0., 0.
*END INSTANCE
*NSET, NSET=LEFT, INSTANCE=ROD_A
1
*NSET, NSET=RIGHT, INSTANCE=ROD_A
2
*END ASSEMBLY

*SUPPORT, SUPPORT_COLLECTOR=BC
LEFT, 0., 0., 0., NAN, NAN, NAN
RIGHT, NAN, 0., 0., NAN, NAN, NAN
*CLOAD, LOAD_COLLECTOR=P
RIGHT, 500., 0., 0., 0., 0., 0.

*LOADCASE, TYPE=LINEARSTATIC
*SUPPORTS
BC
*LOADS
P
*END
\end{femasterdeck}
\begin{abaqusdeck}
*MATERIAL, NAME=STEEL
*ELASTIC
210000., 0.3

*PART, NAME=ROD
*NODE
1,   0., 0., 0.
2, 500., 0., 0.
*ELEMENT, TYPE=T3D2, ELSET=BAR
1, 1, 2
*SOLID SECTION, ELSET=BAR,
 MATERIAL=STEEL
50.
*END PART

*ASSEMBLY, NAME=ASSEMBLY
*INSTANCE, NAME=ROD_A, PART=ROD
500., 0., 0.
*END INSTANCE
*NSET, NSET=LEFT, INSTANCE=ROD_A
1
*NSET, NSET=RIGHT, INSTANCE=ROD_A
2
*END ASSEMBLY

*STEP
*STATIC
*BOUNDARY
LEFT, 1, 3, 0.
RIGHT, 2, 3, 0.
*CLOAD
RIGHT, 1, 500.
*END STEP
\end{abaqusdeck}
\end{dialectexamples}

\section{Homogeneous shell cantilever}

A four-node shell panel is clamped along one edge and loaded transversely at the free edge. The native deck selects \val{MITC4}. Abaqus \val{S4} is mapped by the current reader to the corresponding FEMaster FRT shell implementation. Both Sections represent a homogeneous 2 mm shell; the Abaqus form explicitly uses the supported five Simpson integration points.

\begin{dialectexamples}
\begin{femasterdeck}
*NODE
1,    0.,   0., 0.
2, 1000.,   0., 0.
3, 1000., 500., 0.
4,    0., 500., 0.
*ELEMENT, TYPE=MITC4, ELSET=PANEL
1, 1, 2, 3, 4
*NSET, NAME=ROOT
1, 4
*NSET, NAME=TIP
2, 3

*MATERIAL, NAME=ALUMINIUM
*ELASTIC
70000., 0.33
*SHELL SECTION, ELSET=PANEL,
 MATERIAL=ALUMINIUM, TYPE=INTEGRATED
2.0

*SUPPORT, SUPPORT_COLLECTOR=BC
ROOT, 0., 0., 0., 0., 0., 0.
*CLOAD, LOAD_COLLECTOR=P
TIP, 0., 0., -50., 0., 0., 0.

*LOADCASE, TYPE=LINEARSTATIC
*SUPPORTS
BC
*LOADS
P
*END
\end{femasterdeck}
\begin{abaqusdeck}
*NODE
1,    0.,   0., 0.
2, 1000.,   0., 0.
3, 1000., 500., 0.
4,    0., 500., 0.
*ELEMENT, TYPE=S4, ELSET=PANEL
1, 1, 2, 3, 4
*NSET, NSET=ROOT
1, 4
*NSET, NSET=TIP
2, 3

*MATERIAL, NAME=ALUMINIUM
*ELASTIC
70000., 0.33
*SHELL SECTION, ELSET=PANEL,
 MATERIAL=ALUMINIUM
2.0, 5

*STEP
*STATIC
*BOUNDARY
ROOT, 1, 6, 0.
*CLOAD
TIP, 3, -50.
*END STEP
\end{abaqusdeck}
\end{dialectexamples}

\section{Solid cube under pressure}

The solid example demonstrates an element-face Surface and pressure loading. The native deck uses \cmd{PLOAD}; Abaqus input uses \cmd{DSLOAD} with type \val{P}. In both cases the traction direction follows the finite-element surface orientation rather than an independently supplied global vector.

\begin{dialectexamples}
\begin{femasterdeck}
*NODE
1, 0., 0., 0.
2, 1., 0., 0.
3, 1., 1., 0.
4, 0., 1., 0.
5, 0., 0., 1.
6, 1., 0., 1.
7, 1., 1., 1.
8, 0., 1., 1.
*ELEMENT, TYPE=C3D8, ELSET=CUBE
1, 1, 2, 3, 4, 5, 6, 7, 8
*NSET, NAME=BASE
1, 2, 3, 4
*SURFACE, NAME=LOADED_FACE
CUBE, S2

*MATERIAL, NAME=MAT
*ELASTIC
10000., 0.3
*SOLID SECTION, ELSET=CUBE, MATERIAL=MAT

*SUPPORT, SUPPORT_COLLECTOR=BC
BASE, 0., 0., 0., NAN, NAN, NAN
*PLOAD, LOAD_COLLECTOR=P
LOADED_FACE, 1.0

*LOADCASE, TYPE=LINEARSTATIC
*SUPPORTS
BC
*LOADS
P
*END
\end{femasterdeck}
\begin{abaqusdeck}
*NODE
1, 0., 0., 0.
2, 1., 0., 0.
3, 1., 1., 0.
4, 0., 1., 0.
5, 0., 0., 1.
6, 1., 0., 1.
7, 1., 1., 1.
8, 0., 1., 1.
*ELEMENT, TYPE=C3D8, ELSET=CUBE
1, 1, 2, 3, 4, 5, 6, 7, 8
*NSET, NSET=BASE
1, 2, 3, 4
*SURFACE, NAME=LOADED_FACE,
 TYPE=ELEMENT
CUBE, S2

*MATERIAL, NAME=MAT
*ELASTIC
10000., 0.3
*SOLID SECTION, ELSET=CUBE, MATERIAL=MAT

*STEP
*STATIC
*BOUNDARY
BASE, 1, 3, 0.
*DSLOAD
LOADED_FACE, P, 1.0
*END STEP
\end{abaqusdeck}
\end{dialectexamples}

\section{Eigenfrequency analysis}

The modal example uses a two-element bar with distributed material density and requests the lowest five eigenpairs. The native deck could additionally use \cmd{POINTMASS}; the current Abaqus reader does not translate an equivalent concentrated-mass keyword, so the common example intentionally uses only distributed mass and therefore represents the same mechanical model on both sides.

\begin{dialectexamples}
\begin{femasterdeck}
*NODE
1,    0., 0., 0.
2,  500., 0., 0.
3, 1000., 0., 0.
*ELEMENT, TYPE=T3, ELSET=BAR
1, 1, 2
2, 2, 3
*NSET, NAME=FIXED
1
*NSET, NAME=GUIDED
2, 3
*MATERIAL, NAME=STEEL
*ELASTIC
210000., 0.3
*DENSITY
7.85e-9
*TRUSS SECTION, ELSET=BAR, MATERIAL=STEEL
100.
*SUPPORT, SUPPORT_COLLECTOR=BC
FIXED, 0., 0., 0., NAN, NAN, NAN
GUIDED, NAN, 0., 0., NAN, NAN, NAN

*LOADCASE, TYPE=EIGENFREQ
*SUPPORTS
BC
*NUMEIGENVALUES
5
*END
\end{femasterdeck}
\begin{abaqusdeck}
*NODE
1,    0., 0., 0.
2,  500., 0., 0.
3, 1000., 0., 0.
*ELEMENT, TYPE=T3D2, ELSET=BAR
1, 1, 2
2, 2, 3
*NSET, NSET=FIXED
1
*NSET, NSET=GUIDED
2, 3
*MATERIAL, NAME=STEEL
*ELASTIC
210000., 0.3
*DENSITY
7.85e-9
*SOLID SECTION, ELSET=BAR, MATERIAL=STEEL
100.

*STEP, PERTURBATION
*FREQUENCY
5
*BOUNDARY
FIXED, 1, 3, 0.
GUIDED, 2, 3, 0.
*END STEP
\end{abaqusdeck}
\end{dialectexamples}

\section{Arc-length snap-through}

A shallow two-bar arch illustrates geometrically nonlinear path following. The native deck selects \val{ARC_LENGTH}; Abaqus input selects \cmd{STATIC, RIKS}. The first four Riks data values correspond to the initial, period, minimum, and maximum path scales that the reader normalizes onto the FEMaster nonlinear load path.

\begin{dialectexamples}
\begin{femasterdeck}
*NODE
1, -1000.,   0., 0.
2,     0., 100., 0.
3,  1000.,   0., 0.
*ELEMENT, TYPE=T3, ELSET=BARS
1, 1, 2
2, 2, 3
*NSET, NAME=ENDS
1, 3
*NSET, NAME=APEX
2
*MATERIAL, NAME=STEEL
*ELASTIC
210000., 0.3
*TRUSS SECTION, ELSET=BARS, MATERIAL=STEEL
100.
*SUPPORT, SUPPORT_COLLECTOR=BC
ENDS, 0., 0., 0., NAN, NAN, NAN
APEX, 0., NAN, 0., NAN, NAN, NAN
*CLOAD, LOAD_COLLECTOR=P
APEX, 0., -20000., 0., 0., 0., 0.

*LOADCASE, TYPE=NONLINEARSTATIC
*SUPPORTS
BC
*LOADS
P
*NONLINEAR, CONTROL=ARC_LENGTH,
 MAX_INCREMENTS=200,
 INITIAL_INCREMENT=0.05,
 MINIMUM_INCREMENT=1e-5,
 MAXIMUM_INCREMENT=0.1,
 MAXITER=40, TOL=1e-8
*END
\end{femasterdeck}
\begin{abaqusdeck}
*NODE
1, -1000.,   0., 0.
2,     0., 100., 0.
3,  1000.,   0., 0.
*ELEMENT, TYPE=T3D2, ELSET=BARS
1, 1, 2
2, 2, 3
*NSET, NSET=ENDS
1, 3
*NSET, NSET=APEX
2
*MATERIAL, NAME=STEEL
*ELASTIC
210000., 0.3
*SOLID SECTION, ELSET=BARS, MATERIAL=STEEL
100.

*STEP, NLGEOM=YES, INC=200
*STATIC, RIKS
0.05, 1.0, 1e-5, 0.1
*BOUNDARY
ENDS, 1, 3, 0.
APEX, 1, 1, 0.
APEX, 3, 3, 0.
*CLOAD
APEX, 2, -20000.
*END STEP
\end{abaqusdeck}
\end{dialectexamples}

\section{Linear transient ramp load}

The transient model applies an axial force with a tabular amplitude. Native FEMaster defines the integration interval and Newmark parameters directly. Abaqus \cmd{DYNAMIC, DIRECT} supplies the fixed time increment and Step period; the reader maps these values to the same transient loadcase, whose default Newmark parameters are \(\beta=0.25\) and \(\gamma=0.5\).

\begin{dialectexamples}
\begin{femasterdeck}
*NODE
1,    0., 0., 0.
2, 1000., 0., 0.
*ELEMENT, TYPE=T3, ELSET=BAR
1, 1, 2
*NSET, NAME=FIXED
1
*NSET, NAME=TIP
2
*MATERIAL, NAME=STEEL
*ELASTIC
210000., 0.3
*DENSITY
7.85e-9
*TRUSS SECTION, ELSET=BAR, MATERIAL=STEEL
100.
*SUPPORT, SUPPORT_COLLECTOR=BC
FIXED, 0., 0., 0., NAN, NAN, NAN
TIP, NAN, 0., 0., NAN, NAN, NAN
*AMPLITUDE, NAME=RAMP
0., 0.
0.1, 1.
1., 1.
*CLOAD, LOAD_COLLECTOR=P, AMPLITUDE=RAMP
TIP, 1000., 0., 0., 0., 0., 0.

*LOADCASE, TYPE=LINEARTRANSIENT
*SUPPORTS
BC
*LOADS
P
*TIME
0., 1., 0.01
*NEWMARK
0.25, 0.5
*WRITEEVERY, TYPE=STEPS
10
*END
\end{femasterdeck}
\begin{abaqusdeck}
*NODE
1,    0., 0., 0.
2, 1000., 0., 0.
*ELEMENT, TYPE=T3D2, ELSET=BAR
1, 1, 2
*NSET, NSET=FIXED
1
*NSET, NSET=TIP
2
*MATERIAL, NAME=STEEL
*ELASTIC
210000., 0.3
*DENSITY
7.85e-9
*SOLID SECTION, ELSET=BAR, MATERIAL=STEEL
100.
*AMPLITUDE, NAME=RAMP
0., 0., 0.1, 1., 1., 1.

*STEP
*DYNAMIC, DIRECT
0.01, 1.0
*BOUNDARY
FIXED, 1, 3, 0.
TIP, 2, 3, 0.
*CLOAD, AMPLITUDE=RAMP
TIP, 1, 1000.
*NODE FILE
U
*END STEP
\end{abaqusdeck}
\end{dialectexamples}

\section{Linear harmonic frequency sweep}

The harmonic example evaluates a real axial harmonic force from 10 to 1000 frequency units at 100 linearly spaced points. Native \cmd{FREQUENCIES} constructs the sweep directly. Abaqus \cmd{STEADY STATE DYNAMICS, DIRECT} is translated into the same explicit FEMaster frequency vector.

\begin{dialectexamples}
\begin{femasterdeck}
*NODE
1, 0., 0., 0.
2, 1000., 0., 0.
*ELEMENT, TYPE=T3, ELSET=BAR
1, 1, 2
*NSET, NAME=FIXED
1
*NSET, NAME=TIP
2
*MATERIAL, NAME=STEEL
*ELASTIC
210000., 0.3
*DENSITY
7.85e-9
*TRUSS SECTION, ELSET=BAR, MATERIAL=STEEL
100.
*SUPPORT, SUPPORT_COLLECTOR=BC
FIXED, 0., 0., 0., NAN, NAN, NAN
TIP, NAN, 0., 0., NAN, NAN, NAN
*CLOAD, LOAD_COLLECTOR=HARMONIC
TIP, 1., 0., 0., 0., 0., 0.

*LOADCASE, TYPE=LINEARHARMONIC
*SUPPORTS
BC
*LOADS
HARMONIC
*FREQUENCIES, SCALE=LINEAR
10., 1000., 100
*DAMPING, TYPE=RAYLEIGH
0., 1e-5
*CONSTRAINTMETHOD, TYPE=NULLSPACE
*END
\end{femasterdeck}
\begin{abaqusdeck}
*NODE
1, 0., 0., 0.
2, 1000., 0., 0.
*ELEMENT, TYPE=T3D2, ELSET=BAR
1, 1, 2
*NSET, NSET=FIXED
1
*NSET, NSET=TIP
2
*MATERIAL, NAME=STEEL
*ELASTIC
210000., 0.3
*DENSITY
7.85e-9
*SOLID SECTION, ELSET=BAR, MATERIAL=STEEL
100.

*STEP
*STEADY STATE DYNAMICS, DIRECT,
 FREQUENCYSCALE=LINEAR
10., 1000., 100, 1., 1.
*BOUNDARY
FIXED, 1, 3, 0.
TIP, 2, 3, 0.
*CLOAD, REAL
TIP, 1, 1.
*END STEP
\end{abaqusdeck}
\end{dialectexamples}

\section{Using the examples as templates}

When adapting an example, preserve the modelling hierarchy before tuning numerical details. First verify reference geometry and element connectivity, then define reusable regions, then attach materials and Sections, then define loads and constraints, and only then adjust solver controls. For imported Abaqus input, remain within the keyword subset documented in this manual: a keyword available in Abaqus itself is not automatically implemented by FEMaster's Abaqus reader.

The side-by-side layout is also a translation guide. If a native feature has no Abaqus counterpart in the right column, do not add an assumed Abaqus spelling to an imported deck. Conversely, Abaqus Step grammar should not be copied into a native FEMaster \cmd{LOADCASE}. Both routes converge on the same solver-facing ModelData and loadcase classes while retaining distinct front-end abstractions where that distinction is useful.
